\documentclass[12pt]{article}
\usepackage{amsmath, amssymb}
\usepackage[a4paper, margin=1in]{geometry}

\begin{document}

\section*{�� Lecture 18 Scribe}

\textbf{Marginal PMF, Correlation, Covariance, Joint Gaussian RVs}

\section*{1. Marginal PMF}

\textbf{Definition}

The key idea:

The marginal PMF of one variable is obtained by summing the joint PMF over all possible values of the other variable.

\textbf{Notation}

If joint PMF is $p_{X,Y}(x,y)$, then:

Marginal PMF of $X$:
\[
p_X(x) = \sum_{y} p_{X,Y}(x,y)
\]

Marginal PMF of $Y$:
\[
p_Y(y) = \sum_{x} p_{X,Y}(x,y)
\]

\textbf{Marginal PMF from Joint PMF Table}

Step-by-step procedure:

Start with the joint PMF table

To find $p_X(x)$:
\begin{itemize}
\item Fix $x$
\item Sum across all $y$
\end{itemize}

To find $p_Y(y)$:
\begin{itemize}
\item Fix $y$
\item Sum across all $x$
\end{itemize}

\textbf{Example (as per slides)}

Given: Joint PMF

Solution Steps:

Apply definition:

\[
p_X(x) = \sum_{y} p_{X,Y}(x,y)
\]

Similarly:

\[
p_Y(y) = \sum_{x} p_{X,Y}(x,y)
\]

Compute values by summing corresponding rows/columns

\textbf{Conclusion:}

Marginal PMFs are obtained directly by summation over the joint PMF.

\section*{2. Correlation}

\textbf{Definition}

A basic way to measure joint behavior:

\[
R_{XY} = E[XY]
\]

\textbf{Interpretation}

It gives the average value of product $XY$

Called a raw correlation measure

\textbf{Important Observations}

Depends on scale of X and Y

Computed about the origin (not mean)

Mixes:
\begin{itemize}
\item Mean values
\item Spread
\item Dependence
\end{itemize}

�� Therefore:

Not a clean measure of dependence

\textbf{Special Case}

\[
R_{XY} = 0 \Rightarrow \text{orthogonal about origin}
\]

\textbf{Transition to Covariance}

To remove mean effects:

We center variables $\rightarrow$ leads to covariance

\section*{3. Covariance}

\textbf{Definition}

\[
\text{Cov}(X,Y) = E[(X - E[X])(Y - E[Y])]
\]

\textbf{Step-by-Step Interpretation}

Subtract mean from each variable

Measure how deviations vary together

Meaning:

Positive $\rightarrow$ move in same direction

Negative $\rightarrow$ move in opposite direction

\section*{4. Covariance: Algebraic Form}

Start with:

\[
\text{Cov}(X,Y) = E[(X - E[X])(Y - E[Y])]
\]

Expand:

\[
= E[XY - X E[Y] - Y E[X] + E[X]E[Y]]
\]

Apply linearity of expectation:

\[
= E[XY] - E[X]E[Y] - E[Y]E[X] + E[X]E[Y]
\]

Simplify:

\[
\text{Cov}(X,Y) = E[XY] - E[X]E[Y]
\]

\section*{5. Properties of Covariance}

\textbf{1. Symmetry}

\[
\text{Cov}(X,Y) = \text{Cov}(Y,X)
\]

\textbf{2. Variance as Special Case}

\[
\text{Var}(X) = \text{Cov}(X,X)
\]

\textbf{3. Linearity with Constants}

(Covariance behaves linearly when constants are involved)

\textbf{4. Independence Property}

\[
X \perp Y \Rightarrow \text{Cov}(X,Y) = 0
\]

⚠️ Important:

Zero covariance $\Rightarrow$ uncorrelated

But NOT necessarily independent

\section*{6. Covariance Example 1}

\textbf{Method (as given)}

Use formula:

\[
\text{Cov}(X,Y) = E[XY] - E[X]E[Y]
\]

Substitute:

\[
E[XY]
\]

\[
E[X]
\]

\[
E[Y]
\]

Simplify step-by-step

Final result obtained after substitution

\section*{7. Pearson’s Correlation Coefficient}

\textbf{Definition}

Measures effect of change in one variable when the other changes

\textbf{Key Points}

Measures strength and direction of linear relationship

Normalized version of covariance

Unit-free

\section*{8. Covariance Example 2 (Continuous Case)}

\textbf{Given:}

Joint PDF $f_{X,Y}(x,y)$

\textbf{Steps:}

Compute:

\[
E[X], E[Y], E[XY]
\]

Use:

\[
\text{Cov}(X,Y) = E[XY] - E[X]E[Y]
\]

Perform integration over valid region

Substitute results

Final covariance obtained

\section*{9. Jointly Gaussian Random Variables}

\textbf{Definition}

A pair $(X,Y)$ is jointly Gaussian if:

Joint density has bivariate Gaussian form

Marginals are also Gaussian

\textbf{Interpretation}

Models continuous random variables influenced by many small effects

\textbf{Real-Life Applications}

Signal processing

Communication systems

Noisy measurements

\end{document}